% ============================================================
%  Presentación Beamer — ICS2122 Capstone — Grupo 10
%  Ruteo de Profesionales de la Salud (TD-HDARP)
%  Tiempo objetivo: 7 minutos
%
%  Compilar con: pdflatex presentacion_grupo10.tex
%  Para Fira Sans (metropolis nativo): xelatex presentacion_grupo10.tex
% ============================================================

\documentclass[10pt, aspectratio=169]{beamer}

\usetheme{metropolis}
% \usetheme{Madrid}  % Alternativa si metropolis no está disponible

% ── Paquetes ────────────────────────────────────────────────
\usepackage[T1]{fontenc}
\usepackage[utf8]{inputenc}
\usepackage[spanish]{babel}
\usepackage{booktabs}
\usepackage{amsmath, amssymb}
\usepackage{tikz}
\usepackage{xcolor}
\usepackage{colortbl}
\usepackage{array}

\usetikzlibrary{arrows.meta, positioning}

% ── Paleta de colores ────────────────────────────────────────
\definecolor{ucazul}{RGB}{0, 51, 102}
\definecolor{ucgris}{RGB}{100, 100, 100}
\definecolor{verdeok}{RGB}{0, 128, 64}
\definecolor{rojono}{RGB}{180, 30, 30}
\definecolor{naranja}{RGB}{210, 100, 0}
\definecolor{amarillo}{RGB}{220, 180, 0}

% ── Colores Beamer ───────────────────────────────────────────
\setbeamercolor{frametitle}{bg=ucazul, fg=white}
\setbeamercolor{progress bar}{fg=ucazul}
\setbeamercolor{alerted text}{fg=ucazul}
\setbeamercolor{block title}{bg=ucazul!12, fg=ucazul}
\setbeamercolor{block body}{bg=ucazul!5}
\setbeamercolor{example block title}{bg=verdeok!15, fg=verdeok!70!black}
\setbeamercolor{example block body}{bg=verdeok!5}

\setbeamertemplate{frame numbering}[fraction]

% ── Metadatos ────────────────────────────────────────────────
\title{Ruteo de Profesionales de la Salud}
\subtitle{TD-HDARP · Grupo 10 · ICS2122}
\author{A.~Matta · V.~Encalada · J.A.~Muñoz · N.~Provoste · B.~Suazo}
\institute{Pontificia Universidad Católica de Chile}
\date{Marzo 2026}

% ────────────────────────────────────────────────────────────
\begin{document}

% ============================================================
% PORTADA
% ============================================================
\begin{frame}[plain, noframenumbering]
  \titlepage
  \vspace{-0.8em}
  \begin{columns}[T]
    \column{0.55\textwidth}
      \footnotesize\textbf{Grupo 10}\\[3pt]
      Alonso Matta Sotomayor\\
      Vicente Encalada Henríquez\\
      Juan Andrés Muñoz Solano\\
      Nicolás Provoste Moya \quad Benjamín Suazo Massaro
    \column{0.45\textwidth}
      \footnotesize\textbf{Equipo docente}\\[3pt]
      Profesor: Steffan Heinz\\
      Ayudante: Fernanda Godoy
  \end{columns}
\end{frame}

% ============================================================
% SLIDES 1–2: DESCRIPCIÓN DEL PROBLEMA
% ============================================================
\begin{frame}{Descripción del Problema}

  \begin{block}{El caso}
    Una empresa prestadora de servicios de salud en \textbf{Santiago de Chile} debe planificar diariamente el transporte de sus profesionales a los centros médicos donde trabajan.
  \end{block}

  \vspace{12pt}
  \begin{itemize}\setlength\itemsep{10pt}
    \item Cada día los profesionales solicitan transporte para el día siguiente
    \item Hay dos turnos: \textcolor{verdeok}{\bfseries AM} (casa $\to$ trabajo) y \textcolor{rojono}{\bfseries PM} (trabajo $\to$ casa)
    \item Las rutas se planifican durante la tarde anterior
    \item Puntualidad crítica: impacta directamente la atención de pacientes
  \end{itemize}

\end{frame}

% ─────────────────────────────────────────────────────────────
\begin{frame}{Descripción del Problema \textnormal{\small (cont.)}}

  \begin{itemize}\setlength\itemsep{12pt}
    \item La empresa dispone de una \textbf{flota heterogénea}: vehículos de distintos tamaños y capacidades
    \item Cada vehículo tiene un \textbf{costo fijo} (por usarlo, aunque vaya vacío) y un \textbf{costo variable} (por kilómetro recorrido)
    \item El objetivo es \textbf{minimizar el costo total} cumpliendo estrictamente todas las restricciones
    \item Las restricciones provienen de contratos sindicales, límites de tiempo y características de la flota
  \end{itemize}

  \vspace{14pt}
  \begin{block}{}
    \centering\textbf{TD-HDARP}: Time-Dependent Heterogeneous Dial-a-Ride Problem \quad$\Rightarrow$\quad \textbf{NP-Hard}
  \end{block}

\end{frame}

% ============================================================
% SLIDES 3–6: PRINCIPALES RESTRICCIONES
% ============================================================
\begin{frame}{Principales Restricciones — Categorías Sindicales}

  \begin{block}{La regla}
    Los profesionales se dividen en \textbf{4 categorías} (1 a 4). Un vehículo puede transportar simultáneamente \textbf{a lo más 2 categorías, y deben ser consecutivas}.
  \end{block}

  \vspace{10pt}
  \begin{itemize}\setlength\itemsep{10pt}
    \item \textcolor{verdeok}{\textbf{Permitido:}} categorías 1+2, 2+3, o 3+4 (consecutivas entre sí)
    \item \textcolor{rojono}{\textbf{No permitido:}} categorías 1+3, 1+4, 2+4 (no consecutivas)
    \item Un pasajero \textit{puede} viajar solo si su viaje directo ya supera el tiempo máximo permitido
    \item Es la restricción más inusual del problema — pocos trabajos en la literatura la modelan explícitamente
  \end{itemize}

\end{frame}

% ─────────────────────────────────────────────────────────────
\begin{frame}{Principales Restricciones — Ventanas de Tiempo}

  \begin{block}{La regla}
    Cada profesional tiene una \textbf{ventana de tiempo} que define cuándo puede ser recogido o dejado. El vehículo no puede llegar antes ni después de ese rango.
  \end{block}

  \vspace{10pt}
  \begin{itemize}\setlength\itemsep{10pt}
    \item Turno \textcolor{verdeok}{\bfseries AM}: recogida en domicilio dentro del horario acordado, llegada a tiempo al turno
    \item Turno \textcolor{rojono}{\bfseries PM}: entrega en domicilio dentro del horario post-turno
    \item Las ventanas son \textbf{restricciones duras}: no hay penalización si se violan, simplemente la solución se vuelve infactible
    \item Planificar múltiples rutas respetando todas las ventanas simultáneamente es lo que genera la complejidad combinatoria
  \end{itemize}

\end{frame}

% ─────────────────────────────────────────────────────────────
\begin{frame}{Principales Restricciones — Permanencia en el Vehículo}

  \begin{block}{La regla}
    Existe un \textbf{límite de tiempo} que un pasajero puede permanecer dentro del vehículo durante su traslado.
  \end{block}

  \vspace{10pt}
  \begin{itemize}\setlength\itemsep{10pt}
    \item Evita que un profesional esté en ruta durante mucho tiempo mientras se recogen otros pasajeros
    \item Restricción vinculada al \textbf{bienestar laboral}: forma parte del contrato sindical
    \item Excepción: si el viaje directo entre domicilio y centro médico ya supera el límite, el profesional puede viajar solo en el vehículo
    \item Introduce una tensión entre \textit{eficiencia de la ruta} y \textit{comodidad del pasajero}
  \end{itemize}

\end{frame}

% ─────────────────────────────────────────────────────────────
\begin{frame}{Principales Restricciones — Flota y Estructura de Costos}

  \begin{block}{La regla}
    La empresa opera con vehículos de \textbf{distintos tamaños y capacidades}, cada uno con su propio esquema de costos.
  \end{block}

  \vspace{10pt}
  \begin{itemize}\setlength\itemsep{10pt}
    \item \textbf{Costo fijo}: se incurre al usar el vehículo, independientemente de si va lleno o vacío
    \item \textbf{Costo variable}: proporcional a los kilómetros recorridos
    \item La decisión de qué vehículo asignar a cada ruta afecta directamente el costo total
    \item Usar muchos vehículos pequeños o pocos grandes son estrategias con trade-offs distintos
  \end{itemize}

\end{frame}

% ============================================================
% SLIDE 7: ANÁLISIS DE DATOS (EN DESARROLLO)
% ============================================================
\begin{frame}{Análisis de Datos}

  \vspace{20pt}
  \begin{block}{}
    \centering
    \vspace{8pt}
    \textit{Esta sección se encuentra actualmente en desarrollo.}\\[4pt]
    Los datos del caso de estudio están siendo recopilados y procesados.\\[4pt]
    Se presentará en una entrega futura.
    \vspace{8pt}
  \end{block}

\end{frame}

% ============================================================
% SLIDE 8: KPIs
% ============================================================
\begin{frame}{KPIs — Indicadores de Desempeño}

  \vspace{4pt}
  Los siguientes indicadores permitirán evaluar la calidad de cualquier solución propuesta:

  \vspace{10pt}
  \begin{columns}[T]
    \column{0.50\textwidth}
      \begin{itemize}\setlength\itemsep{10pt}
        \item \textbf{Costo operacional total}\\
          \footnotesize Suma de costos fijos + variables de todos los vehículos utilizados
        \item \textbf{Tiempo promedio de permanencia}\\
          \footnotesize Tiempo medio que cada pasajero pasa dentro del vehículo. Aunque el máximo es restricción dura, el promedio diferencia la calidad entre soluciones factibles
        \item \textbf{Tasa de ocupación de la flota}\\
          \footnotesize Pasajeros transportados sobre la capacidad total despachada
      \end{itemize}

    \column{0.50\textwidth}
      \begin{itemize}\setlength\itemsep{10pt}
        \item \textbf{Distancia total recorrida}\\
          \footnotesize Kilómetros acumulados por toda la flota en operación
        \item \textbf{Número de vehículos utilizados}\\
          \footnotesize Cada vehículo adicional implica un costo fijo; minimizar flota es un objetivo secundario relevante
        \item \textbf{Tiempo de cómputo}\\
          \footnotesize Tiempo real de ejecución del algoritmo. El contexto operativo exige una solución en menos de 60 minutos para ser usable en producción
      \end{itemize}

      \vspace{10pt}
      \begin{block}{}
        \footnotesize Serán evaluados comparando la solución propuesta contra una línea base heurística.
      \end{block}
  \end{columns}

\end{frame}

% ============================================================
% SLIDES 9–10: REVISIÓN BIBLIOGRÁFICA / ESTADO DEL ARTE
% ============================================================
\begin{frame}{Estado del Arte — El DARP en la literatura}

  \begin{itemize}\setlength\itemsep{10pt}
    \item El \textbf{Dial-a-Ride Problem (DARP)} surgió para modelar transporte de personas con movilidad reducida y pacientes — su lógica se extendió hoy al transporte corporativo y de salud
    \item \textbf{Ho et al.\ (2018)} revisaron 277 publicaciones y encontraron que el área ha crecido exponencialmente en los últimos 20 años, impulsada por aplicaciones reales en múltiples industrias
    \item \textbf{Cordeau \& Laporte (2003)} demostraron que el DARP es \textbf{NP-Hard}: ningún algoritmo exacto puede resolverlo eficientemente para instancias reales
    \item Nuestro caso (\textbf{TD-HDARP}) combina usuarios heterogéneos + tiempos dependientes de la hora: una variante con menos de 20 papers específicos en la literatura
  \end{itemize}

\end{frame}

% ─────────────────────────────────────────────────────────────
\begin{frame}{Estado del Arte — Hallazgos Clave}

  \begin{itemize}\setlength\itemsep{10pt}
    \item \textbf{Cordeau (2006)} demostró que los métodos exactos resuelven óptimamente solo hasta \textasciitilde{}16 solicitudes; con más pasajeros, el tiempo de cómputo supera las horas
    \item \textbf{Ropke \& Pisinger (2006)} revolucionaron el campo: con \textbf{ALNS} manejaron 100+ solicitudes manteniendo alta calidad de solución — hoy es el estado del arte
    \item \textbf{Parragh (2011)} encontró que la \textbf{heterogeneidad de usuarios} es el factor más complejo; requiere operadores especializados que respeten las reglas de compatibilidad
    \item \textbf{Ichoua et al.\ (2003)} mostraron que ignorar los tiempos de viaje dependientes de la hora genera rutas infactibles en la práctica — la congestión no puede omitirse
  \end{itemize}

\end{frame}

% ============================================================
% SLIDES 11–13: DISCUSIÓN METODOLÓGICA
% ============================================================
\begin{frame}{Metodología — Métodos Exactos (MILP)}

  \begin{block}{¿En qué consiste?}
    Modelar el problema como un \textbf{programa lineal entero mixto} y resolverlo con algoritmos Branch-and-Cut o Branch-and-Price (usando solucionadores como Gurobi).
  \end{block}

  \vspace{8pt}
  \begin{columns}[T]
    \column{0.50\textwidth}
      \begin{exampleblock}{\textcolor{verdeok}{Ventajas}}
        \begin{itemize}\setlength\itemsep{5pt}
          \item Garantiza la \textbf{solución óptima} global
          \item Proporciona una \textbf{cota inferior} para comparar otros algoritmos
          \item Útil para validar soluciones en instancias pequeñas
        \end{itemize}
      \end{exampleblock}

    \column{0.50\textwidth}
      \begin{alertblock}{\textcolor{rojono}{Limitaciones}}
        \begin{itemize}\setlength\itemsep{5pt}
          \item Complejidad \textbf{exponencial} (NP-Hard)
          \item Inviable para más de \textasciitilde{}40 solicitudes: supera 1 hora de cómputo
          \item Nuestro caso tiene \textbf{cientos de profesionales} → completamente impracticable
        \end{itemize}
      \end{alertblock}
  \end{columns}

\end{frame}

% ─────────────────────────────────────────────────────────────
\begin{frame}{Metodología — Heurísticas Constructivas}

  \begin{block}{¿En qué consiste?}
    Construir una solución \textbf{paso a paso} asignando pasajeros a rutas de forma secuencial, eligiendo en cada paso la opción ``más barata'' disponible (\textit{inserción voraz / greedy}).
  \end{block}

  \vspace{8pt}
  \begin{columns}[T]
    \column{0.50\textwidth}
      \begin{exampleblock}{\textcolor{verdeok}{Ventajas}}
        \begin{itemize}\setlength\itemsep{5pt}
          \item \textbf{Extremadamente rápidas}: cientos de pasajeros en segundos
          \item Fáciles de implementar y de entender
          \item Útiles como \textbf{solución inicial} para otros algoritmos
        \end{itemize}
      \end{exampleblock}

    \column{0.50\textwidth}
      \begin{alertblock}{\textcolor{rojono}{Limitaciones}}
        \begin{itemize}\setlength\itemsep{5pt}
          \item Visión \textbf{cortoplacista}: decisiones locales llevan a rutas globalmente malas
          \item Quedan atrapadas en \textbf{óptimos locales} de baja calidad
          \item La restricción de categorías sindicales genera frecuentes \textbf{rutas infactibles}
        \end{itemize}
      \end{alertblock}
  \end{columns}

\end{frame}

% ─────────────────────────────────────────────────────────────
\begin{frame}{Metodología — Metaheurísticas (ALNS)}

  \begin{block}{¿En qué consiste?}
    El \textbf{Adaptive Large Neighborhood Search} destruye y reconstruye iterativamente partes de la solución, usando operadores especializados que se adaptan según su desempeño previo.
  \end{block}

  \vspace{8pt}
  \begin{columns}[T]
    \column{0.50\textwidth}
      \begin{exampleblock}{\textcolor{verdeok}{Ventajas}}
        \begin{itemize}\setlength\itemsep{5pt}
          \item Escala a instancias \textbf{reales} con cientos de solicitudes
          \item Permite diseñar operadores que \textbf{respetan la regla de categorías}
          \item Es el \textbf{estado del arte} en problemas de ruteo complejo (Ropke \& Pisinger, 2006)
          \item El Recocido Simulado integrado permite escapar de óptimos locales
        \end{itemize}
      \end{exampleblock}

    \column{0.50\textwidth}
      \begin{alertblock}{\textcolor{rojono}{Limitaciones}}
        \begin{itemize}\setlength\itemsep{5pt}
          \item Implementación más \textbf{compleja} que las heurísticas simples
          \item Requiere \textbf{ajuste cuidadoso de parámetros} para converger bien
          \item No garantiza el óptimo global
        \end{itemize}
      \end{alertblock}
  \end{columns}

\end{frame}

% ============================================================
% SLIDE 14: METODOLOGÍA SELECCIONADA
% ============================================================
\begin{frame}{Metodología Seleccionada — ALNS}

  \begin{block}{¿Por qué ALNS?}
    Es la única metodología que \textbf{escala al tamaño real del problema} y permite manejar la complejidad de la restricción sindical de categorías.
  \end{block}

  \vspace{8pt}
  \begin{itemize}\setlength\itemsep{9pt}
    \item La restricción de categorías \textit{rompe} el espacio de búsqueda: las heurísticas simples generan rutas infactibles de forma sistemática
    \item El ALNS puede diseñar un operador de destrucción específico que \textbf{identifica y repara} los conflictos de categoría en cada iteración
    \item Parragh (2011) demostró que este enfoque funciona en problemas análogos con usuarios heterogéneos
    \item Permite obtener soluciones de alta calidad dentro del \textbf{límite operativo de 60 minutos}
  \end{itemize}

\end{frame}

% ============================================================
% SLIDE 15: PASOS FUTUROS (EN DESARROLLO)
% ============================================================
\begin{frame}{Pasos Futuros}

  \vspace{20pt}
  \begin{block}{}
    \centering
    \vspace{8pt}
    \textit{Esta sección se encuentra actualmente en desarrollo.}\\[4pt]
    El plan de trabajo detallado se presentará en una entrega futura.
    \vspace{8pt}
  \end{block}

\end{frame}

% ============================================================
% DECLARACIÓN DE USO DE IA (sin numeración)
% ============================================================
\begin{frame}[plain, noframenumbering]{Declaración de Uso de IA}

  \begin{columns}[T]
    \column{0.50\textwidth}
      \begin{block}{Herramientas utilizadas}
        \footnotesize
        \vspace{4pt}
        \textbf{Gemini:} Redacción, ortografía y búsqueda bibliográfica\\[4pt]
        \textbf{Claude:} Asistencia en código \LaTeX\ y análisis
        \vspace{4pt}
      \end{block}
    \column{0.50\textwidth}
      \begin{alertblock}{Consideración del equipo}
        \footnotesize
        El análisis, decisiones metodológicas y argumentos presentados son responsabilidad exclusiva del equipo.
      \end{alertblock}
  \end{columns}

  \vspace{14pt}
  \centering\small\textcolor{ucgris}{Grupo 10 · ICS2122 · Pontificia Universidad Católica de Chile · Marzo 2026}

\end{frame}

% ────────────────────────────────────────────────────────────
\end{document}
